\section*{Erd\H{o}s Problem 172: simultaneous finite sums and products}
\subsection*{Statement and outcome}
For a finite $A\subset\mathbb N$, let $\operatorname{FS}(A)$ and
$\operatorname{FP}(A)$ consist respectively of sums and products over nonempty subsets, with no repeated elements. The question asks whether every finite colouring admits arbitrarily large $A$ for which their union has one colour.

We prove the assertion for eventually periodic colourings and for a class of simultaneous prime-valuation colourings. For arbitrary colourings we prove a canonical reduction to a finite list of colours indexed by subset size. Making that list one colour remains unresolved in this attempt.

\subsection*{Eventually periodic colourings}
Suppose $c(n+q)=c(n)$ for all $n\ge N_0$. All sufficiently large multiples of $q$ then have a common colour. Choose any desired number of distinct multiples of $q$ at least $\max(N_0,2)$. Every nonempty sum and product is again such a multiple, so the whole union has this colour. This includes the simpler monochromatic-ideal observation in the old local attempt.

\subsection*{Simultaneous valuation colourings}
Let $P$ be a finite nonempty set of primes, and let the colour of $n$ depend only on the tuple
$(v_p(n)\bmod m_p)_{p\in P}$, for fixed positive integers $m_p$.
Put $M=\operatorname{lcm}_{p\in P}m_p$, $B=\prod_{p\in P}p$, and choose
\[
 A=\{B^{M},B^{2M},\ldots,B^{rM}\}.
\]
For a nonempty subset $S$ of the exponents, its sum factors as
$B^{Mj_0}(1+\sum_{j\in S,\ j>j_0}B^{M(j-j_0)})$,
where $j_0=\min S$. The parenthesis is one modulo every $p\in P$.
Thus each $p$-valuation of this sum is exactly $Mj_0$.
The valuation of the corresponding product is $M\sum_{j\in S}j$.
All these valuations are zero modulo their respective $m_p$.
Hence all sums and products have the colour assigned to the zero tuple.
Distinct valuations are essential here: unrestricted sums of terms of equal valuation could introduce cancellation.

\subsection*{A reduction valid for every finite colouring}
Fix a target size $r$. There is an infinite set $B\subset\mathbb N$ and colours
$\gamma_j,\delta_j$, $1\le j\le r$, such that every $j$-subset $S\subset B$ satisfies
\[
 c\left(\sum_{s\in S}s\right)=\gamma_j,\qquad
 c\left(\prod_{s\in S}s\right)=\delta_j.
 \tag{1}
\]
Apply the infinite Ramsey theorem successively for $j=1,\ldots,r$, colouring each $j$-subset by the ordered pair of displayed colours and passing to an infinite homogeneous subset. Later restrictions preserve earlier properties.

Here the needed Ramsey theorem itself follows by induction on uniformity. The case of singletons is infinite pigeonhole. For a colouring of $j$-subsets, recursively choose vertices $v_i$ and nested infinite tails so that the colour of $\{v_i\}\cup T$, for $(j-1)$-subsets $T$ of the new tail, is a fixed colour $\eta_i$; the induction hypothesis supplies that tail. Infinitely many $\eta_i$ agree. The corresponding $v_i$ form a homogeneous set, since the least-indexed member determines the colour of every $j$-subset. This proves the input and hence (1).

Any $r$ distinct members of $B$ now have sums and products whose colours depend only on cardinality. But there are up to $2r-1$ independent entries in this profile, since only $\gamma_1=\delta_1$ is automatic. Ramsey homogeneity does not assert that these entries agree. This is the exact gap in trying to turn (1) into the desired conclusion.

\subsection*{Why clearing rational denominators is insufficient}
If a rational set $A$ has common denominator dividing $D$, replacing it by $DA$ changes a subset sum to $D\sum A'$ but a product of $j$ elements to $D^j\prod A'$. An arbitrary integer colouring need not preserve colours under any of these different dilations. Thus the rational analogue reported on the problem page does not transfer merely by clearing denominators.

\subsection*{Assessment and attribution}
The old \texttt{gpt\_pro\_5.2/172.tex} handled parity and a monochromatic ideal. Those are included in the proved periodic special case; the valuation construction and cardinality-profile reduction give further explicit partial results. Source and literature context: \url{https://www.erdosproblems.com/172}. No external sums-and-products theorem is invoked as a proof, and no novelty claim is made.
\subsection*{COMPLETION ESTIMATE}
\textbf{COMPLETION: 30\%}
